\section{Operación sostenida: la jornada no es una entrega repetida}
\label{sec:operacion}

Todo lo anterior optimiza \emph{una} misión: formar la coalición, certificar el
\emph{wrench}, encaminar y entregar. Una flota industrial no hace una misión:
encadena transporte, espera y recarga durante horas. Esta sección muestra que
ambos objetivos discrepan, y da el marco en el que la discrepancia se mide.

\subsection{Medidas de ocupación con duración}

Sea $z$ el estado operativo de un AMR —carga transportada, batería, situación—
y $b$ una \emph{continuación} admisible desde $z$, de duración $\tau_{zb}$,
recompensa $r_{zb}$ y consumo $e_{zb}$. Sea $q_{zb}\ge0$ la frecuencia por
unidad de tiempo con que se ejecuta esa continuación. El régimen estacionario
satisface el balance de flujo y una normalización que \emph{pesa por duración},
porque no todas las maniobras duran lo mismo:
\begin{equation}
  \sum_b q_{zb}=\sum_{y,b}P(z\mid y,b)\,q_{yb},
  \qquad
  \sum_{z,b}\tau_{zb}\,q_{zb}=1,
  \qquad q_{zb}\ge0 .
  \label{eq:ocupacion}
\end{equation}

\begin{proposicion}[Caudal como programa lineal]
\label{pr:caudal}
Sobre el conjunto \eqref{eq:ocupacion}, el caudal de entregas por unidad de
tiempo es $\sum_{z,b}r_{zb}q_{zb}$ y su máximo es el valor de un programa
lineal. Las restricciones de recurso —energía media, ocupación de cargadores,
disponibilidad de AMR— son lineales en $q$ y se añaden sin cambiar la clase de
problema. El valor óptimo acota superiormente el caudal sostenido de cualquier
política estacionaria admisible.
\end{proposicion}

\begin{proof}
El balance y la normalización son lineales, el objetivo también, y toda política
estacionaria induce una medida de ocupación admisible; el máximo del programa no
puede ser inferior al caudal de ninguna de ellas.
\end{proof}

La formulación anterior es la de un proceso semi-Markov con recompensas, en el
que el tiempo entre transiciones es una variable aleatoria y no una constante
\cite{lozovanu2024positional}; la normalización por duración de
\eqref{eq:ocupacion} es lo que distingue ese caso del markoviano ordinario. Esa
literatura obliga además a declarar una hipótesis que el enunciado anterior
omite.

\begin{supuesto}[Cadena única]
\label{s:unichain}
La cadena inducida por toda política estacionaria admisible es \emph{unichain}:
su matriz límite tiene todas las filas iguales, o equivalentemente, las
probabilidades límite no dependen del estado inicial.
\end{supuesto}

\begin{observacion}[Qué falla sin esa hipótesis]
\label{obs:multichain}
Sin \cref{s:unichain} el valor de \cref{pr:caudal} sigue siendo una cota sobre el
conjunto de medidas de ocupación admisibles, pero deja de ser alcanzable desde un
estado dado. El contraejemplo es mínimo: con tres estados en que el primero cae
irremediablemente en una clase recurrente de recompensa nula y el tercero forma
otra clase de recompensa unitaria, el programa lineal devuelve $1{,}0000$
—colocando toda la masa en la clase buena— mientras que la recompensa media
alcanzable es $0{,}0000$ desde los dos primeros estados y $1{,}0000$ solo desde
el tercero. La matriz límite tiene dos filas distintas, que es la firma del caso
\emph{multichain}.

Para una flota real la hipótesis se lee así: todo estado degradado ha de ser
recuperable, de modo que exista una única clase recurrente de operación. Si un
AMR varado fuera irrecuperable, el modelo se partiría en clases y el
dimensionamiento basado en \cref{pr:caudal} prometería un caudal que la
instalación no puede dar.
\end{observacion}

\begin{observacion}[El rescate entra en el caudal como un número]
\label{obs:rescate-caudal}
La conexión con \cref{sec:standby} no es solo cualitativa. Modelando la flota con
tres situaciones —operativa, averiada y recargando— y haciendo variar la tasa de
rescate $p$, el valor del programa lineal es $0$ para $p=0$, $0{,}273$ para
$p=0{,}02$, $0{,}600$ para $p=0{,}10$ y $0{,}774$ para $p=0{,}40$, y en todos los
casos con $p>0$ coincide con lo alcanzable desde cualquier estado inicial hasta
la cuarta cifra.

Queda así cerrada una cadena que atraviesa tres secciones: el emplazamiento de la
reserva determina la cobertura de \cref{def:cobertura}; la cobertura \emph{es} la
función $p$ de \cref{th:subprovision}; y $p$ entra en este programa lineal
fijando el caudal. Nivel de reserva, dónde ponerla y cuánto se produce son tres
lecturas de un mismo número.

Conviene señalar también qué ocurre en el extremo: con $p=0$ el programa devuelve
exactamente cero, y lo hace correctamente, porque la avería pasa a ser la única
clase recurrente. El modelo no falla ahí; avisa.
\end{observacion}

\subsection{Óptimo por entrega frente a óptimo por jornada}

\begin{observacion}[Una coalición rápida puede ser mala para la jornada]
\label{obs:eco-fast}
Con batería en $\{0,1,2,3\}$ y tres continuaciones —\textsf{ECO}: $5$~s, consume
$1$, entrega $1$; \textsf{FAST}: $3$~s, consume $3$, entrega $1$; y recarga:
$4$~s, repone $1$— el ciclo puro \textsf{FAST} dura $3+3\times4=15$~s y da
$0{,}0667$ entregas/s, mientras el ciclo \textsf{ECO} dura $5+4=9$~s y da
$0{,}1111$. El programa lineal \eqref{eq:ocupacion} devuelve exactamente
$1/9=0{,}1111$.

\textsf{FAST} gana la entrega aislada —tres segundos frente a cinco— y pierde la
jornada por un factor $1{,}67$. La selección de coalición que solo mira la
misión en curso está optimizando la magnitud equivocada.
\end{observacion}

\begin{observacion}[El transitorio no persiste]
\label{obs:transitorio}
Una enumeración de políticas deterministas arrancando con batería llena da
$0{,}1121$, ligeramente por encima del valor estacionario. La diferencia es la
ventaja de partir cargado, que se consume una sola vez. El valor del programa
lineal es el caudal de largo plazo, y es el que debe usarse para dimensionar.
\end{observacion}

Esa diferencia tiene nombre y estructura propia, y conviene dárselos, porque
resulta ser el criterio que falta cuando el caudal no decide.

\begin{proposicion}[Ganancia y sesgo]
\label{pr:sesgo}
Para una política estacionaria con cadena aperiódica, el valor acumulado en
horizonte $T$ desde el estado $z$ se descompone como
\[
  V_T(z)\;=\;g\,T\;+\;h(z)\;+\;o(1),
\]
donde $g$ es la ganancia —el caudal de \cref{pr:caudal}— y $h$ el \emph{sesgo},
solución de $(I-P+\Pi)h=r-g\mathbf 1$ normalizada por $\Pi h=0$. El sesgo mide la
ventaja o desventaja de partir de un estado concreto, y se consume una sola vez
\cite{lozovanu2024positional}.
\end{proposicion}

\begin{observacion}[El caudal no siempre decide, y entonces decide el sesgo]
\label{obs:sesgo-desempate}
Dos cadenas con la misma ganancia pueden tener sesgos muy distintos. Con dos
estados y recompensa unitaria en el primero, la cadena de mezcla rápida
—$P_{ij}=1/2$— y la de mezcla lenta —permanencia $0{,}9$— tienen ambas
$g=0{,}5$, y sesgos $(0{,}5;-0{,}5)$ y $(2{,}5;-2{,}5)$ respectivamente: un factor
cinco. La descomposición se verifica a precisión de máquina, con error
$10^{-13}$ para $T=2000$.

La consecuencia para el dimensionamiento es concreta. \Cref{pr:caudal} puede
tener varias soluciones óptimas con idéntico caudal, y el programa lineal no
distingue entre ellas. El sesgo sí: prefiere las políticas cuya ventaja de
partida se agota antes, es decir, las de mezcla rápida. Y esa preferencia
coincide con la que impone la capa de comunicación, porque la mezcla rápida es
también lo que acorta la complejidad temporal de \cref{def:complejidad}. Cuando
el caudal empata, conviene romper el empate por sesgo, y hacerlo no exige
información nueva: se calcula del mismo modelo.

Lo medido en \cref{obs:transitorio} es, por tanto, un sesgo.
\end{observacion}

\begin{observacion}[Media no es garantía]
\label{obs:media-energia}
Una restricción de energía \emph{media} en \eqref{eq:ocupacion} no impide agotar
la batería en una realización concreta. La seguridad energética se conserva
admitiendo únicamente continuaciones compatibles con el estado real y su
reserva, es decir, restringiendo el conjunto admisible por estado, no el
promedio. Es el mismo principio que separa $\hat\theta$ de $\Theta$ en
\cref{obs:adaptacion}.
\end{observacion}

\subsection{El puente exacto entre ocupación, precios de estado y revisión}

\Cref{pr:caudal} plantea la jornada como programa de ocupación con criterio
promedio. Para una misión que \emph{termina} —hay entrega— la formulación
pertinente es la absorbente, y su dual es lo que enlaza esta capa con el residual
de \cref{th:bellman-residual} \cite{mayorga2026m1}.

\begin{proposicion}[Ocupación absorbente, precios de estado y coste reducido]
\label{pr:ocupacion-absorbente}
Sea $m_{zb}\ge0$ el número esperado de ejecuciones de la continuación $b$ desde
el estado $z$ antes de la entrega. Los balances en los estados no terminales son
$\sum_b m_{zb}-\sum_{y,b}P(z\mid y,b)\,m_{yb}=\nu_z$, y el problema de ocupación
es $\min\{c^{\!\top}m:\ Bm=\nu,\ m\ge0\}$. Su dual es
\[
  \max_V\ \nu^{\!\top}V
  \quad\text{s.a.}\quad
  V(z)-\sum_{z'}P(z'\mid z,b)\,V(z')\ \le\ c(z,b),
\]
y el coste reducido $A_V(z,b)=c(z,b)+\sum_{z'}P(z'\mid z,b)V(z')-V(z)$ es no
negativo, nulo sobre el soporte de la ocupación óptima, y coincide con el
residual de Bellman de \cref{th:bellman-residual} tomando $W=V$.
\end{proposicion}

\begin{observacion}[Comprobación, y qué cuenta $m$]
\label{obs:ocupacion-num}
Sobre un proceso absorbente de cinco estados y tres continuaciones por estado,
el primal y el dual coinciden a $1{,}3\times10^{-15}$; el coste reducido mínimo
es exactamente cero y su valor absoluto sobre el soporte de $m$ no supera
$4{,}4\times10^{-16}$. La política recuperada normalizando $m_{zb}$ por estado
coincide con la de $\arg\min_b A_V(z,b)$ en los cinco estados visitados.

Dos precisiones que la derivación exige. Primera: $m$ cuenta usos esperados;
\emph{no se reparte físicamente una carga entre futuros incompatibles}, y cada
realización ejecuta una continuación atómica y segura. La ocupación es una
herramienta de valoración, no una fracción de la carga. Segunda: la política
estacionaria se recupera normalizando $m_{zb}$ por su suma en el estado, con
tratamiento aparte de los estados no visitados, donde $m$ no informa.
\end{observacion}

\begin{observacion}[Los precios de estado dan el incentivo correcto]
\label{obs:precios-estado}
Esto cierra un círculo que estaba abierto entre tres secciones. El pago que la
dinámica de revisión debe ver para elegir continuaciones es $-A_V$: con él, la
comparación por pares prefiere las continuaciones de coste reducido nulo, que son
las óptimas, y \cref{obs:equipo-barato} deja de ser una advertencia para
convertirse en regla —comparar $c+PV-V$, no $c$—. Los precios de estado $V$ son
un dual como lo son los precios de \emph{wrench} y de corredor, y se difunden
del mismo modo; y el residual $\varrho$ que \cref{th:bellman-residual} tolera por
unidad de duración es el error con que se conoce $A_V$, que
\cref{pr:presupuesto-error} acota. Lo que sigue sin cubrir es la identificación
de $c$ y $P$ para una instalación concreta, que es la calibración que
\cref{obs:operacion-abierto} declara pendiente.
\end{observacion}

\subsection{Un solo potencial para seleccionar y para ejecutar}

Hay una exigencia de coherencia que atraviesa \cref{th:potencial},
\cref{th:testigo} y \cref{obs:precios-estado} y que conviene enunciar como
principio: \emph{la función que selecciona participantes debe valorar el mismo
problema que autoriza ejecutar} \cite{mayorga2026m1}. Una formulación coherente
—todavía no equivalente al catálogo aproximado de la rama física, y así se
declara— es el valor de misión
\[
  V_t(d)=\inf_{(X,U)\in\Tcal(d,Y_t)}\Phi_t(d,X,U),
  \qquad
  \Phi_t=V_f(X_{t+H})+\int_t^{t+H}\bigl[\ell_0\mathbf 1_{X\notin G}+\ell_E+\ell_U+\ell_R+\ell_{\mathrm{cambio}}\bigr]ds,
\]
con unidades normalizadas por escalas declaradas y pesos no negativos. Tres
rasgos lo distinguen de un coste de trayectoria ordinario: se cuenta la energía
de \emph{toda} la flota relevante, incluido el AMR excluido que debe retirarse;
el conjunto objetivo $G$ exige posición, orientación, velocidades y
\emph{permanencia}, y la descarga vertical es un requisito aparte; y el término
$\ell_{\mathrm{cambio}}$ cobra las reconfiguraciones, que es lo que hace del
autómata de \cref{def:automata} un objeto valorado y no solo secuenciado. Con
$J_i=\Phi$ para todos —juego de interés común— o con costes marginales cuya
diferencia unilateral sea la del potencial, el pseudogradiente
$F=\operatorname{col}_i\nabla_{z_i}J_i$ es el de \cref{th:vi}, y el reposo de la
revisión y el óptimo del testigo son el mismo punto.

\begin{proposicion}[Batería sin regeneración y reserva en línea]
\label{pr:bateria}
Sin recuperación de energía, una aproximación de la potencia de batería del
AMR $i$ es
\[
  P^{\mathrm{bat}}_i=P^{\mathrm{aux}}_i+P^{\mathrm{loss}}_i
  +\sum_{s\in\{L,R\}}\frac{[\tau_{is}\,\omega_{w,is}]_+}{\eta_{is}},
  \qquad
  \dot E_i=-P^{\mathrm{bat}}_i,
  \qquad
  E_i\ \ge\ E^{\mathrm{retorno}}_i+E^{\mathrm{reserva}}_i ,
\]
y la desigualdad debe comprobarse en línea sobre la realización, no solo
prefijarse en el catálogo.
\end{proposicion}

\begin{observacion}[Lo que el corchete positivo cambia]
\label{obs:bateria-num}
La potencia mecánica negativa —frenado, descenso— no se convierte
automáticamente en recarga. Sobre una misión de $1500$~s con un paquete de
$100$~Wh, contar esa potencia como recarga deja $90{,}2$~Wh; con el corchete
positivo quedan $61{,}5$~Wh. La diferencia, $28{,}7$~Wh, es un $29\%$ del paquete:
un catálogo que trate el frenado como recuperación sobreestima la autonomía en
esa magnitud y autoriza continuaciones que la reserva no cubre. Es la versión
con número de \cref{obs:media-energia}: la restricción media de \cref{pr:caudal}
dimensiona la jornada, pero la desigualdad de reserva se comprueba en cada
realización con el estado real, y una reserva prefijada en el catálogo no
sustituye esa comprobación.
\end{observacion}

\subsection{Programar la recarga: la ranura ficticia}

\Cref{obs:media-energia} advierte que una restricción de energía media no protege
una realización concreta, y \cref{tab:posicionamiento} declara la gestión de
recarga como abordada solo parcialmente. Hay una construcción que la cubre, y su
interés está en un detalle de formulación.

El obstáculo es que un presupuesto energético es una \emph{desigualdad} —cada AMR
necesita al menos $E_k$— mientras que la maquinaria de \cref{obs:analogia} exige
un simplex, es decir, una igualdad. El puente es añadir una \emph{ranura
ficticia} que absorbe la holgura \cite{martinezPiazuelo2025gne}: con ella el
presupuesto pasa de «$\le$» a «$=$» sin cambiar el problema, y toda la
invariancia queda disponible.

\begin{observacion}[Un programa de recarga con tres restricciones a la vez]
\label{obs:recarga}
Con cuatro AMR, cinco franjas horarias, tarifa con valle en la tercera, tope de
potencia por AMR y capacidad del cuadro $C=2{,}2$, el mecanismo produce un
programa que satisface simultáneamente tres cosas de tres maneras distintas:
\begin{enumerate}[label=(\roman*),leftmargin=2.2em,topsep=0.2em,itemsep=0.15em]
  \item la masa del simplex se conserva exactamente —$6{,}000000$, $5{,}000000$,
  $7{,}500000$, $5{,}000000$—, por la estructura de la dinámica;
  \item el tope de potencia por AMR y franja se respeta en todo instante, porque
  se impone de forma prescriptiva, multiplicando la tasa por el margen restante
  (\cref{obs:prescriptivo});
  \item la capacidad del cuadro se excede en $2{,}211$ frente a $2{,}2$ —un
  $0{,}5\%$—, porque se impone con precio, y un precio solo actúa en el
  equilibrio.
\end{enumerate}
Las tres restricciones del mismo problema se comportan exactamente como
\cref{obs:criterio-restricciones} predice según el instrumento con que se
imponen. Es la comprobación más directa de ese criterio que contiene el trabajo.
\end{observacion}

\begin{observacion}[Qué hace el mecanismo cuando no hay energía para todos]
\label{obs:racionamiento}
En el ensayo la demanda total, $11{,}5$, supera a la capacidad agregada,
$5\times2{,}2=11$. El mecanismo no falla ni declara infactibilidad: el precio de
capacidad sube hasta $8{,}7$–$9{,}6$, próximo a la penalización de no cargar,
y el déficit se reparte —$0{,}17$, $0$, $0{,}18$ y $0{,}30$ entre los cuatro
AMR—. Que el precio sombra de la capacidad tienda al valor de la demanda no
servida es lo económicamente correcto, y significa que el racionamiento emerge
del mismo mecanismo en lugar de requerir una lógica aparte.

Lo que sigue sin abordarse es el \emph{emplazamiento} de las estaciones, que es
una decisión estratégica distinta. La programación temporal, en cambio, queda
cubierta por esta construcción.
\end{observacion}

\subsection{La capa discreta como juego posicional}

La selección dentro del catálogo de continuaciones es un problema sobre un grafo
de estados y decisiones, que es exactamente la clase de los \emph{juegos
posicionales} \cite{lozovanu2024positional}: los jugadores controlan posiciones
de un grafo dirigido, las estrategias estacionarias puras asignan a cada
posición controlada un sucesor, y el interés está en la existencia de
equilibrios en estrategias \emph{estacionarias y puras} —no aleatorizadas—,
que es lo que un ejecutor determinista puede implementar.

\begin{observacion}[Qué aporta esa lectura]
\label{obs:posicional}
Dos cosas concretas. Primera, cuando el grafo de dependencias entre cargas es
acíclico —por disjunción de equipos o por precedencia de recursos—, la selección
óptima se calcula por suma mínima sobre el grafo en tiempo polinómico, sin
enumerar el catálogo. Segunda, la existencia de equilibrio en estrategias
estacionarias puras es lo que legitima que el mecanismo devuelva una decisión
determinista y reproducible, en lugar de una distribución sobre continuaciones.
Ambas son propiedades de la capa discreta y no se transfieren a las variables
continuas de \cref{def:estrategia}.

La segunda requiere una hipótesis. La
existencia de equilibrio de Nash en estrategias estacionarias \emph{puras} está
demostrada para los juegos posicionales estocásticos con criterio descontado, y
con criterio promedio \emph{bajo la propiedad de cadena única}
\cite{lozovanu2024positional}. Es la misma \cref{s:unichain} que exige
\cref{pr:caudal}, y no es casualidad: ambas afirmaciones dependen de que el
comportamiento asintótico no dependa del estado de partida. Sin ella, lo que se
garantiza es equilibrio en estrategias mixtas, y el mecanismo tendría que
devolver una distribución sobre continuaciones en lugar de una decisión.
\end{observacion}

\begin{observacion}[Y hay una segunda condición, sobre lo que se observa]
\label{obs:pura-observacion}
La garantía de estrategia pura tiene además un supuesto tácito que conviene
sacar a la luz: se refiere al juego sobre el \emph{estado}. Un AMR no decide
sobre el estado sino sobre lo que observa a $n$ saltos, y dos estados distintos
pueden producir la misma observación. En ese régimen —el de las cadenas
parcialmente observables— la política óptima sobre observaciones es en general
\emph{aleatorizada} \cite{clempner2024optimization}.

La comprobación es tajante. Con tres estados de los que dos son indistinguibles
y una acción distinta óptima en cada uno, las cuatro políticas deterministas
sobre observaciones rinden $0{,}22439$, mientras que aleatorizar al observar la
clase ambigua rinde $0{,}40690$: un $81\%$ más. Una política determinista debe
tratar igual a dos estados que exigen acciones opuestas; aleatorizar reparte el
acierto entre ambos.

La consecuencia para este trabajo es de alcance, no de corrección: la
reproducibilidad que \cref{obs:posicional} atribuye al mecanismo vale mientras la
información local baste para distinguir las situaciones que exigen decisiones
distintas. Cuando no baste —dos coaliciones candidatas indistinguibles a $n$
saltos— insistir en determinismo tiene un coste medible, y la alternativa
correcta es aleatorizar. Es otra vía por la que la calidad de la información
entra en el rendimiento, junto a \cref{obs:tensado-coste}.
\end{observacion}

\subsection{Las estadísticas de la jornada son exactas, no asintóticas}

Los indicadores de la capa discreta —fracción de flota en reserva, caudal
esperado, probabilidad de cumplir un plazo— se han tratado hasta aquí como
magnitudes a simular. No hace falta: tienen forma cerrada.

\begin{teorema}[Distribución estacionaria de Gibbs]
\label{th:gibbs}
Sea $\Psi$ el potencial exacto de \cref{th:potencial} sobre el conjunto finito
de configuraciones $\sigma$, y considérese la revisión asíncrona que propone un
cambio según la regla logit y lo acepta con probabilidad de Metropolis. Entonces
la cadena es reversible y su distribución estacionaria es exactamente
\[
  \pi(\sigma)=\frac{e^{-\beta\Psi(\sigma)}}{Z},
  \qquad Z=\sum_{\sigma'}e^{-\beta\Psi(\sigma')} .
\]
En consecuencia: (i) la distribución es exacta para $N$ finito, sin invocar
población grande; (ii) todo indicador expresable como esperanza bajo $\pi$ se
obtiene de $Z$ y sus derivadas; (iii) al crecer $\beta$, $\pi$ se concentra en
$\arg\min\Psi$.
\end{teorema}

\begin{proof}
Con propuesta $q$ y aceptación
$\alpha(\sigma\!\to\!\sigma')=\min\{1,\,q(\sigma|\sigma')e^{-\beta\Psi(\sigma')}/
(q(\sigma'|\sigma)e^{-\beta\Psi(\sigma)})\}$ se verifica el balance detallado
$\pi(\sigma)P(\sigma,\sigma')=\pi(\sigma')P(\sigma',\sigma)$ por construcción,
luego $\pi$ es estacionaria; la irreducibilidad sobre el catálogo da unicidad.
(iii) es inmediato de la forma exponencial.
\end{proof}

\begin{observacion}[Comprobación]
\label{obs:gibbs-num}
Sobre una cadena de seis configuraciones con potencial aleatorio y $\beta=2$, la
estacionaria calculada como vector propio del operador de transición coincide con
$e^{-\beta\Psi}/Z$ con error $1{,}1\times10^{-16}$: exactitud numérica, no
aproximación.
\end{observacion}

\begin{observacion}[El mismo recorrido calcula el óptimo y la partición]
\label{obs:semianillo}
El cálculo de $\min_\sigma\Psi$ por programación dinámica sobre el catálogo y el
cálculo de $Z$ son el mismo recorrido del mismo grafo con dos pares de
operadores: $(\min,+)$ devuelve el óptimo y $(+,\times)$ devuelve la partición.
El puente entre ambos es
\[
  -\beta^{-1}\ln\!\sum_{y}e^{-\beta h(y)}
  \;=\;\min_y h(y)\;-\;\beta^{-1}\ln\!\sum_{y}e^{-\beta(h(y)-\min h)},
  \qquad\text{luego}\quad
  0\le \min_y h(y)+\beta^{-1}\ln\!\sum_y e^{-\beta h(y)}\le\frac{\ln|Y|}{\beta}.
\]
Con $|Y|=12$ el error medido es $0{,}018$ para $\beta=20$ frente a una cota de
$0{,}124$, y $1{,}8\times10^{-4}$ para $\beta=100$ frente a $0{,}025$.

Es la misma desigualdad de \emph{log-sum-exp} que acota la brecha de la
envolvente compuesta en \cref{pr:barrera-compuesta}. Aparece dos veces por la
misma razón: en ambos sitios se sustituye un extremo por una versión suave para
poder derivar, y en ambos el precio de la suavización es logarítmico en el número
de términos.
\end{observacion}

\begin{observacion}[Lo que cuesta la exactitud]
\label{obs:gibbs-precio}
Dos precios, y conviene decirlos. El primero: la exactitud de \cref{th:gibbs}
exige el test de aceptación. Con propuesta logit pura la estacionaria es de Gibbs
solo si los vecindarios de revisión son regulares —si todas las configuraciones
ofrecen el mismo conjunto de alternativas—, condición que un catálogo real no
cumple. Añadir la aceptación de Metropolis es barato y convierte una
estacionaria aproximada en exacta.

El segundo: $\beta$ grande concentra en el óptimo pero enlentece la exploración,
porque el tiempo de escape de una cuenca de profundidad $\Delta$ crece como
$e^{\beta\Delta}$. Precisión de la selección y velocidad de exploración son el
mismo parámetro tirando en sentidos opuestos, y esa tensión no se elimina: se
elige dónde situarse en ella.
\end{observacion}

\subsection{El reposo del logit no es un equilibrio de Nash}

Hay un tercer precio de la temperatura, más serio que los dos anteriores porque
afecta a \emph{dónde} se detiene el mecanismo y no solo a cuándo.

\begin{proposicion}[Estacionariedad de Nash perturbada]
\label{pr:nash-perturbado}
Las dinámicas evolutivas de comparación por pares —como la de Smith, que gobierna
el reparto de esfuerzo— y las de exceso de pago separable son \emph{Nash
estacionarias}: se detienen exactamente en $\arg\max_{\tilde x}\tilde x^{\!\top}p$.
Las de mejor respuesta perturbada —el logit, que gobierna el reclutamiento— no lo
son. Satisfacen en su lugar la estacionariedad de Nash \emph{perturbada}: se
detienen en $\arg\max_{\tilde x}\{\tilde x^{\!\top}p-\omega(\tilde x)\}$, donde
$\omega$ es la perturbación asociada al protocolo \cite{martinezPiazuelo2025gne}.
\end{proposicion}

\begin{observacion}[Cuánto se desplaza, medido]
\label{obs:desplazamiento-logit}
La consecuencia es que el reposo del reclutamiento no es el equilibrio del juego
sino el del juego perturbado, y la distancia entre ambos es $\Theta(1/\beta)$.
Sobre un juego de población con tres estrategias, el producto
$\beta\lVert x_{\mathrm{logit}}-x_{\mathrm{Nash}}\rVert$ converge a
$0{,}275$ para un equilibrio interior —$0{,}228$ con $\beta=10$ hasta $0{,}275$
con $\beta=200$— y se estabiliza en torno a $0{,}37$ cuando el equilibrio está en
la frontera. Es decir, la desviación decae como el inverso de la temperatura
inversa, sin anularse nunca a $\beta$ finito.

La temperatura tiene por tanto \emph{tres} efectos y no dos: concentra la
distribución estacionaria (\cref{th:gibbs}), enlentece la mezcla de forma
exponencial (\cref{obs:gibbs-precio}) y desplaza el reposo respecto del
equilibrio exacto. Los tres empujan en direcciones incompatibles: subir $\beta$
corrige el desplazamiento y arruina la exploración.
\end{observacion}

\begin{observacion}[Qué obliga a corregir en los enunciados]
\label{obs:nash-perturbado-alcance}
Conviene precisar el alcance de lo afirmado en otras secciones. Lo demostrado
sobre el reparto de esfuerzo y los precios no se ve afectado, porque esas capas
usan comparación por pares y son Nash estacionarias. Lo que queda condicionado es
la lectura del reposo del \emph{reclutamiento}: no debe leerse como un equilibrio
del juego original, sino del perturbado, con la desviación acotada arriba.

Esto es coherente con lo que \cref{obs:ruido} ya sostenía —el logit se emplea
como perturbación de exploración y no como descripción del reposo— pero lo hace
cuantitativo: la diferencia entre ambas lecturas es $\Theta(1/\beta)$ y no
cualitativa.
\end{observacion}

\begin{observacion}[Los AMR pueden decidir con reglas distintas]
\label{obs:sociedad-mixta}
Un corolario práctico compensa lo anterior. Una sociedad \emph{mixta}, en la que
unas poblaciones revisan por mejor respuesta perturbada y otras por comparación
por pares o exceso de pago, sigue satisfaciendo la condición que sostiene la
estabilidad asintótica —en su versión perturbada—
\cite{martinezPiazuelo2025gne}. La arquitectura de este trabajo \emph{es} una
sociedad mixta: logit en la capa discreta, comparación por pares en la continua.

La lectura de diseño es más amplia que eso: los AMR no necesitan compartir regla
de decisión. Una flota que mezcle generaciones o fabricantes, con lógicas de
revisión distintas, conserva las garantías mientras cada regla pertenezca a una
de las tres clases. Es una forma de heterogeneidad distinta de la de capacidades —hasta aquí solo eran heterogéneas las capacidades— y que resulta
admisible sin cambiar nada.
\end{observacion}

\subsection{El límite de población}

\Cref{obs:alcance-poblacional} declaró que la abstracción poblacional es
aproximada en el reclutamiento porque $N$ es pequeño y las identidades importan.
El otro extremo tiene su propia teoría.

\begin{observacion}[Qué ocurre cuando la flota crece]
\label{obs:campo-medio}
Al crecer $N$ con la razón cargas/AMR fija, la interacción de cada agente con el
resto pasa a depender solo de la \emph{distribución} de estrategias y no de las
identidades: es el régimen de campo medio \cite{tembine2021meanfield}. En ese
límite la hipótesis de influencia individual despreciable se recupera y el
análisis poblacional se vuelve exacto, con brecha que decrece con $N$. La
consecuencia útil no es el límite en sí, sino su lectura sobre el diseño: la
arquitectura no se degrada de forma catastrófica al crecer la flota —mejora—, y
los resultados atómicos de \cref{th:potencial} y \cref{th:descomposicion} siguen
valiendo en todo el rango porque no invocan población grande. El régimen difícil
es el intermedio, que es donde opera un almacén real.
\end{observacion}

\subsection{Posicionamiento frente a la práctica industrial}

La literatura de automatización de almacenes define el bloqueo como la
incapacidad de varios robots de avanzar porque cada uno pretende ocupar la
posición que otro ocupa \cite{yildirim2023warehouse}: una espera circular. Las
estrategias documentadas son mayoritariamente \emph{reactivas} —detectar el
ciclo, pausar robots, poner el tiempo de tarea a infinito, apagar unos segundos,
o declarar al vecino como obstáculo para forzar el recálculo— y la propia obra
señala que la práctica puntera se desplaza hacia enfoques \emph{proactivos}, que
buscan ruta alternativa antes de que el conflicto ocurra.

\begin{observacion}[Dónde encaja el mecanismo de este trabajo]
\label{obs:posicionamiento-almacen}
El mecanismo de reserva vecinal de este trabajo es proactivo por construcción: la coalición solicita el
testigo \emph{antes} de entrar en la zona compartida, y la terna lexicográfica
$(w_i,P_i,-i)$ decide sin recálculo global. La diferencia con las estrategias
tabuladas en esa literatura no es el mecanismo —las reservas por ventana
temporal existen— sino la garantía: el resultado de ausencia condicional de
inanición acota la espera bajo hipótesis explícitas, mientras que las soluciones
industriales reportadas son heurísticas evaluadas empíricamente. Y el peaje
multiclase de \cref{th:multiclase} añade algo que esa literatura no contempla:
discriminar por la envolvente de la coalición, no solo por prioridad de tarea.
\end{observacion}

\begin{observacion}[Lo que la capa operativa aún no cierra]
\label{obs:operacion-abierto}
\Cref{pr:caudal} da el caudal máximo de un modelo de estados agregado; no
predice el caudal de una instancia concreta con geometría, congestión y fallos.
Enlazar ambos exige poblar $\tau_{zb}$ y $P(z\mid y,b)$ con medidas de campaña, y
esa calibración no se hace aquí. Se declara, por tanto, como marco de
dimensionamiento y no como predicción validada.
\end{observacion}

\subsection{Colas de tareas: estabilidad y lo que la física no deja estabilizar}

La jornada es una cola: las tareas llegan mientras otras se ejecutan. El
criterio de estabilidad es el clásico de servicio por peso máximo
\cite{tassiulas1992stability}, y lo que importa aquí es que se conserva con
servicio aproximado y que tiene un límite físico \cite{mayorga2026compendio}.

\begin{teorema}[Estabilidad con servicio de peso máximo aproximado]
\label{th:maxweight}
Sea $Q_k(t+1)=[Q_k(t)-S_k(t)]_++A_k(t)$ con
$\mathbb E[A(t)\,|\,Q(t)]=\lambda$, segundos momentos uniformemente acotados
y servicios en una región compacta $\Sigma$. Si existe $\epsilon>0$ con
$\lambda+\epsilon\mathbf 1\in\mathrm{conv}(\Sigma)$ y la política satisface
$Q(t)^{\!\top}\mathbb E[S(t)\,|\,Q(t)]\ge\max_{s\in\Sigma}Q(t)^{\!\top}s-B$,
entonces las colas son fuertemente estables y
\begin{equation}
  \label{eq:maxweight}
  \limsup_{T}\frac1T\sum_{t=0}^{T-1}\sum_k\mathbb E[Q_k(t)]\le\frac{B_0+B}{\epsilon},
\end{equation}
con $B_0$ dependiente de los segundos momentos. Si en el modelo reducido el
servicio maximiza $Q^{\!\top}S(s)-G(s)$ con $0\le G\le G_{\max}$, la hipótesis
se cumple con $B=G_{\max}$.
\end{teorema}

\begin{proof}
Con $V=\lVert Q\rVert^2/2$ y $([q-s]_++a)^2\le q^2+s^2+a^2+2q(a-s)$, la
deriva es $\le B_0+Q^{\!\top}\lambda-Q^{\!\top}\mathbb E[S\,|\,Q]\le B_0+B-\epsilon\lVert Q\rVert_1$
porque $\max_sQ^{\!\top}s\ge Q^{\!\top}(\lambda+\epsilon\mathbf 1)$;
telescopar. Para el servicio penalizado,
$Q^{\!\top}S(\hat s)\ge Q^{\!\top}S(s^\star)-G_{\max}$.
\end{proof}

\begin{proposicion}[Imposibilidad fuera de la región física]
\label{pr:imposibilidad-fisica}
Sea $\Sigma_{\mathrm{phys}}$ la región de vectores de servicio que satisfacen
el certificado físico. Si existe $w\ge0$ con
$w^{\!\top}\lambda>\sup_{s\in\Sigma_{\mathrm{phys}}}w^{\!\top}s$, ninguna
política que respete el certificado estabiliza todas las colas con tasa
media de llegadas $\lambda$.
\end{proposicion}

\begin{proof}
El servicio acumulado ponderado en $T$ pasos es a lo sumo
$T\sup_sw^{\!\top}s$, y las llegadas ponderadas crecen como
$Tw^{\!\top}\lambda$; la diferencia, que acota inferiormente
$w^{\!\top}Q(T)$, diverge linealmente.
\end{proof}

\begin{observacion}[El caudal es una región, no un número]
\label{obs:region-caudal}
\Cref{th:maxweight} y \cref{pr:imposibilidad-fisica} dan al programa de
ocupación de \cref{s:unichain} su lectura de colas: el caudal sostenible no
es un escalar sino la envolvente convexa de los vectores de servicio que el
certificado admite, y una demanda fuera de ella no se estabiliza con ninguna
regla de prioridad. El término de retraso acumulado del potencial de misión
es exactamente el peso $Q$ de \cref{th:maxweight}, de modo que negociar con
él es servir por peso máximo salvo una penalización acotada. Con tres colas,
tasas $0{,}3$ y cinco vectores de servicio, una réplica reducida con
penalización uniforme en $[0,0{,}4]$ mantiene el atraso medio en $5{,}1$
tareas durante doscientos mil pasos.
\end{observacion}
